\documentclass{inprogress}
\newcommand{\tr}{\operatorname{tr}}
\title{A Complex valued CLT type argument}
\begin{document}
\maketitle

We start with the following bilinear CLT estimate:

\begin{lemma}
    Let $\mu$ be a  $d_0$-dimensional random variable in $L^2$, with covariance $\mathbb X^{i,j}$.
    Let $N_t \to \infty $ be a sequence natural numbers as $t\in \mathbb N \to \infty$, and $X_0, \dots, X_{N_t}$ be a sequence of IID random variables $\sim \mu_0 \in L^2$, with covariance $\mathbb X^{i,j}$.
    Let $\alpha_{1}, \dots, \alpha_{N_t}$ be a sequence of vectors in $\mathbb R^{d_1}$ satisfying:
    \[\lim_{t\to \infty}\max_{j \in 1, \dots N_t} \|\alpha_t\| \to 0, \quad\text{and} \quad \sum_{s=1}^{N_t} a_s^{(k)}a_s^{(l)} \to \mathbb A^{k,l}\].

    Let $B_{s,t}^u:\mathbb R^{d_0}\times \mathbb R^{d_1} \to \mathbb R^{d}$ be a bilinear form. 

    Let $S_t^u:= \sum_{s=0}^{N_t} a^{(i)}_{s}X^{(j)}_{s} B_{i,j}^u $ (with Einstein notation convention). Then $S_t^{u}$ converges as $t\to \infty$ uniformly to a Normal with mean zero and covariance matrix
    \[
    \Sigma^{u,v} :=   \mathbb X^{i,j} A^{k,l} B_{i,k}^u B_{j,l}^v 
    \]
\end{lemma}

\begin{proof}
    By polarization, it suffices to consider the case $d_0 = d_1 = d = 1$, after wihch it follows from the standard CLT. In fact, polarization gives access to the covariance matrix at finite $t$: Let $\mathbb A_t := \sum_{k=s}^{N_t} a_s^{(k)}a_s^{(l)}$, then 
    \[
    (\operatorname{Cov}\mathbb S_t)^{u,v} :=  \mathbb X^{i,j} A^{k,l}_t B_{i,k}^u B_{j,l}^v \]
\end{proof}

In our application, $d_0 = d_1 = d = 2$, and $B$ is the complex product in $\mathbb C \sim \mathbb R^2$. In this case, we can lower-bound the covariance matrix $\Sigma$.


\begin{lemma} If $d_0=d_1=d = 2$, and $B$ is the complex product in $\mathbb C$,
\[
\det \Sigma \gtrsim \tr(\mathbb A)^2 \cdot \det \mathbb X
\]
\end{lemma}

\begin{proof}
    Note that the determinant is super-additive in the space of Symmetric matrices of dimension 2 (By Minkowski's determinant inequality), and in particular it suffices to check the case when $\mathbb A$ is rank one. By complex rotation symmetry, it suffices to check the case 
    \[
    A = \begin{pmatrix}
        1 & 0 \\ 0 & 0
    \end{pmatrix}.
    \]

    In this case $\tr \mathbb A = 1$ and $\Sigma = \mathbb X$, and the result follows. 
\end{proof}

\begin{lemma}[Minkowski's determinant inequality.] Let $A,B$ be symmetric nonnegative definite $n\times n$ matrices, then
$$
\det A + \det B
\le
\left((\det A)^{1/n}
+
(\det B)^{1/n}\right)^{n}
\le 
\det(A+B)
$$
\end{lemma}

In the concentration-compactness type argument, this shows that if we start with a random variable $\mathbb X$ that is not (a rotation of) a real random variable, then the linear combinations converge to a Gaussian that is quantitatively far from being one-dimensional. 

\section{Trying to understand the SPR question}

We will prove that if $X$ is a $L^2$ random variable that does not take only two real variables and $G$ is a standard Gaussian random variable, then the span of infinitely many copies of $X$ and two copies of $G$ does SPY.

The first realization is that having two copies of $G$ or infinitely many is free. (The only thing that matters is the covariance between the two Gaussian elements we are testing, and that can be simulated with two elements).

The second realization is that CLT tells us that the space of pairs of linear combinations is (in law) closed, and in particular stable spr follows from spr.
\end{document}